% !TeX root = ../../../main.tex
\subsection{Phenotyping Prompt}\label{ch06:apd:phenotype-prompt}
For Sepsis-3 determination, we had a clinician review the discharge summaries and assign each case as meeting the Sepsis-3 definition.  
To develop a computational approach, we constructed a prompt specific to Sepsis-3, which describes the qSOFA procedure for determining a Sepsis-3 diagnosis, and used it as query to Llama-3.1-8B-Instruct (L3.1 8B) and QwQ-32B-Preview (Qwen).  
We used the phenotyping prompt:

\begin{tcolorbox}[breakable,
  colback=yellow!10,
  colframe=orange!80!black,
  title=Phenotyping Prompt,
  boxrule=2pt
]
    {You are an expert physician.}  
    
    {Determine if the patient described in the following case report has either sepsis or septic shock, as defined by the Sepsis-3 criteria, which correspond to having a} 
    
    {(1) suspected or confirmed infection and} 
    
    {(2) blood pressure/respiratory rate/mental status abnormalities.} 
    
    {If the information is not present, use your best judgment based on the information available.} 
    
    {Reply 1 for sepsis, 0 otherwise.} 
    
    {Reply with the number 0 or 1 only in \textbackslash\!boxed\{\textbackslash\!n \boxed{<0 $ or $ 1>} \textbackslash\!n\}with no explanation.}
    
    {Here is the case:} 
    
    {\textless \textbackslash end prefix\textgreater}
\end{tcolorbox}



